% Seção: F-EFI · Eficiência da rede (RF-EFI-01 e RF-EFI-02)
% Resultados experimentais na rede real pendentes; ver feature-F-EFI.md.
\subsection{Eficiência da rede}\label{sec:F-EFI}

A eficiência estrutural é calculada em $G_P$, cuja unidade é a parada. Para cada origem,
uma busca Dijkstra calcula as distâncias aos destinos; a eficiência global é a média de
$1/d(i,j)$ sobre os pares ordenados distintos. Pares sem caminho contribuem zero:
\[
  E(G_P) = \frac{1}{n(n-1)}
    \sum_{i \ne j,\ d(i,j) < \infty} \frac{1}{d(i,j)}.
\]
Como os pesos representam segundos, a unidade do resultado é $\mathrm{s}^{-1}$.
Quando a métrica é amostrada, sorteiam-se origens uniformemente sem reposição; o erro
padrão usa a variância amostral das contribuições e correção de população finita.

O tempo médio de viagem usa $G_E$, que soma caminhada, espera esperada e viagem de ônibus,
como no contrato de F-ROTA. A amostra contém pares ordenados distintos de paradas, não
plataformas; portanto, as arestas de desembarque de peso zero não entram na média de
$1/d$. Uma busca por origem fornece os tempos aos destinos da amostra. Pares sem trajeto
são informados separadamente e não são tratados como viagens de duração zero.

\pendente{Figura da página Eficiência.}

\pendente{Medição com o GTFS real: média e desvio de cinco execuções, tamanho da amostra,
$V$, $E$, máquina e comando. O feed real não estava disponível durante esta implementação.}

\textbf{Limitações.} $G_P$ usa tempos medianos de viagem e caminhada, sem espera. Em $G_E$,
a espera continua no valor provisório de 300~s até a integração da tabela da F-ESPERA. O
tempo médio estimado é condicional aos pares alcançáveis na amostra.
